% ==============================================================================
% ПАРТИЯ B03-015: DISS-005-B0657 .. DISS-005-B0697
% Глава 1: Измерение тока с применением катушек Роговского и ёмкостных делителей. РНМ нулевой последовательности
% Замечания: DISS-005-C0066 .. DISS-005-C0069
% Иллюстрации: image41.jpeg, image42.png .. image47.png
% ==============================================================================

% DISS-005-B0657
\subsection{Измерение тока с применением катушек Роговского и ёмкостных делителей}
\label{sec:ch1_rogowski_and_capacitive}

% DISS-005-B0658
Электромагнитные трансформаторы напряжения имеют достаточную точность преобразования и не вызывают проблем при функционировании БАВР. Внедрение новых типов преобразователей напряжения, таких как ёмкостные [125,126] или резистивные делители [127,128], связано в первую очередь с решением проблем возникновения феррорезонансных явлений, энергоэффективности и массогабаритных показателей.

% DISS-005-B0659
Совершенно иным образом обстоят дела с электромагнитными трансформаторами тока (ТТ), которые могут иметь существенные погрешности преобразования в переходных режимах коротких замыканий при насыщении их магнитопроводов [129-131]. Другими недостатками электромагнитных ТТ являются: сложность установки на существующих подстанциях (в виду высоких массогабаритных показателей и неразборности конструкции), необходимость прокладки медных кабелей, высокую стоимость. В виду описанных недостатков электромагнитных ТТ некоторые производителей устройств релейной защиты и автоматики начали применять катушки Роговского (КР). Данные первичные преобразователи (рис. 2) лишены ряда недостатков электромагнитных ТТ, и могут иметь разъемную конструкцию и при этом обеспечивают достаточную точностью преобразования в широком диапазоне измеряемых токов. КР применяются в комплектных распределительных устройствах (КРУ), например, в КРУ Etalon производства ООО «ТАВРИДА ЭЛЕКТРИК» [122], в шкафах БАВР производства ООО «АПС» [3,4], а также множестве зарубежных установок [132-136]. Следует отметить, что современные устройства РЗА позволяют использовать сигнал с КР как напрямую, например, алгоритмы, работающие на производных тока [119-121], так и производить интегрирование сигнала на основе мгновенных значений или в частотной области, для дальнейшего использования первичного тока в классических алгоритмах.

% DISS-005-B0660
Однако типы датчиков, основанные на катушках Роговского не имеют широкой практики применения в отечественной промышленности, и лишь постепенно приобретает всё более обширное применение. А в устройствах БАВР, до 2023 года вообще отсутствовал. Поэтому данный вопрос требует подробного изучения для оценки возможных улучшений и сложностей, которые возникнут при их применении.

% DISS-005-B0662 .. DISS-005-B0668 (DISS-005-FIG-038, FIG-039, FIG-040 / DISS-005-TAB-026)
\begin{figure}[htbp]
    \centering
    \begin{minipage}{0.32\textwidth}
        \centering
        \includegraphics[width=\textwidth]{image41.jpeg}
        \centerline{\small а)}
    \end{minipage}\hfill
    \begin{minipage}{0.32\textwidth}
        \centering
        \includegraphics[width=\textwidth]{image42.png}
        \centerline{\small б)}
    \end{minipage}\hfill
    \begin{minipage}{0.32\textwidth}
        \centering
        \includegraphics[width=\textwidth]{image43.png}
        \centerline{\small в)}
    \end{minipage}
    \caption{Катушка Роговского «NRC-150»: а) внешний вид; б) установка в ячейке КРУ; в) применения КР в КРУ Etalon производства ООО «ТАВРИДА ЭЛЕКТРИК» (3)}
    \label{fig:ch1_img41_43}
\end{figure}

% DISS-005-B0671
\subsubsection{Информация о применении}

% DISS-005-B0672
С развитием систем БАВР они получаю всё более широкое распространение, причём не только в сетях 6-35~кВ, но и в сетях 0,4~кВ. Например, в июне \commentnote{DISS-005-C0066}{Евдаков Алексей}{Саратов}{2023 }года в опытно промышленную эксплуатацию был установлен соответствующий шкаф БАВР. В указанных сетях, классические трансформаторы тока практически отсутствуют, поэтому необходимо применение иных датчиков тока, которые легко устанавливаются, а именно -- катушек Роговского. Данный шкаф, в декабре 2023 года успешно завершил этап опытно-промышленной эксплуатации и введён в работу. А в ближайшие годы, ввиду успешности проведённых испытаний, будет установлено множество аналогичных устройств на других объектах для повышения качества электроснабжения.

% DISS-005-B0673
Помимо этого, в связи с внедрением систем БАВР совместно с КРУ Etalon производства ООО «ТАВРИДА ЭЛЕКТРИК» посредством применения измерительных преобразовательных модулей (ИПМ) серии МИР-ИПМ, вопрос обработки сигналов с катушек Роговского, а также и ёмкостных делителей становится актуален не только для сетей 0,4~кВ, но и для всех остальных классов напряжения, где применятся данный устройства. Данные устройства так же были внедрены в эксплуатации и успешно функционирую.

% DISS-005-B0674
\subsubsection{Разбор \commentnote{DISS-005-C0067}{Евдаков Алексей}{Уточнение про 0,4кВ}{осциллограмм }и обсуждение вопросов}

% DISS-005-B0675
Ввиду новизны применения данных типов датчиков для устройств БАВР, считаем актуальным вопрос изучения поведения органа РНМ в переходных режимах, которые характерны для БАВР.

% DISS-005-B0676
Осциллограммы срабатывания БАВР и последующего срабатывания ВНР приведены на рисунках 1 и 2. Как можем видеть, для сигналов, снимаемых с КР характерна, излишняя зашумлённость и наличие импульсных бросков в моменты переключения. При этом, особого внимания заслуживает режим параллельной работы секций шин при функционировании ВНР (рис. 2). В данный временной промежуток времени, возникают дополнительные высшие гармоники.

% DISS-005-B0677 .. DISS-005-B0678 (DISS-005-FIG-041 / DISS-005-TAB-027)
\begin{figure}[htbp]
    \centering
    \includegraphics[width=0.85\textwidth]{image44.png}
    \caption{Рисунок 1. Срабатывание БАВР по 1СШ после установки блока бесперебойного питания.}
    \label{fig:ch1_img44}
\end{figure}

% DISS-005-B0680 .. DISS-005-B0681 (DISS-005-FIG-042 / DISS-005-TAB-028)
\begin{figure}[htbp]
    \centering
    \includegraphics[width=0.85\textwidth]{image45.png}
    \caption{Рисунок 2. Срабатывание ВНР 1СШ.}
    \label{fig:ch1_img45}
\end{figure}

% DISS-005-B0684
\subsection{Необходимость проработки РНМ по 0-ой последовательности \commentnote{DISS-005-C0068}{Евдаков Алексей}{Можно оставить задачу, даже если не будет расписано внедрение 67N в наш терминал и его применение.\newline\newlineК тому же это позволит указать актуальность изучение нейронок для анализа ОЗЗ на данных Галины.}{для }\commentnote{DISS-005-C0069}{Евдаков Алексей}{Данные файлы находятся в папке\newline«\_вспомогательные файлы\textbackslash{}Нижнекамск Шина»\newline}{БАВР}}
\label{sec:ch1_zero_sequence_rnm}

% DISS-005-B0685
\subsubsection{Информация о применении и актуальности проработки темы}

% DISS-005-B0686
Ранее ОЗЗ регистрировалось по средствам органа контроля напряжения нулевой последовательности (59N) и выполняло исключительно блокирующие функции.

% DISS-005-B0687
С производств поступают запросы о внедрении в систему БАВР цепей пуска по факту регистрации ОЗЗ. Так как ранее указанная функция отсутствовала, то в классические БАВР не заводились цепи от ТТНП и вообще отсутствовали какие-либо реализации РНМ по нулевой последовательности.

% DISS-005-B0688
К тому же, в ряде устройств БАВР блокировка осуществляется по условию ОЗЗ на обеих секциях, так как данный режим не считается аварийным. И, по сути, защита осуществляется от возможного возникновения двухфазного КЗ ввиду включения СВ и ОЗЗ в разных фазах. Однако в ряде сетей, ОЗЗ является частым событием, и перерастает в двухфазные замыкания, которое может быть внутренним (Рис. 1 и 2).

% DISS-005-B0689
А в случае предварительного срабатывания БАВР (до 2-ух фазного КЗ) по каким-либо причинам, мы подвергаем опасности аварийного режима, по сути, всю подстанцию.

% DISS-005-B0690
Следовательно, вопрос добавление органов, осуществляющих локализацию места ОЗЗ является очень важным. И требуется дополнительное исследование данного направления.

% DISS-005-B0691 .. DISS-005-B0692 (DISS-005-FIG-043 / DISS-005-TAB-029)
\begin{figure}[htbp]
    \centering
    \includegraphics[width=0.85\textwidth]{image46.png}
    \caption{Рисунок 1. Напряжение с ТН СШ с обоих терминалов БАВР за весь промежуток аварии.}
    \label{fig:ch1_img46}
\end{figure}

% DISS-005-B0694 .. DISS-005-B0695 (DISS-005-FIG-044 / DISS-005-TAB-030)
\begin{figure}[htbp]
    \centering
    \includegraphics[width=0.85\textwidth]{image47.png}
    \caption{Рисунок 2. Осциллограмма с терминала РЗА отходящей линии.}
    \label{fig:ch1_img47}
\end{figure}
